%% kapa: exam=2 round=27 year=2016 period=2 subject=appraisal_theory subject_ko=감정평가이론 seq=2 number=2 points=30
%% source: Q-Net 감정평가사 기출문제 (cst00309) artlSeq=5185401 file=extracted/1008/5185401/제27회 감정평가사 2차시험 기출문제/제27회 감정평가사 2차 시험문제지/(2교시 시험문제지)감정평가이론.pdf
%% conversion: kordoc 4.24.51 -> tools/gen2.py (auto); review=proofread-gpt-sol; flags=none
\begin{longquestion}{【문제 2】}
감정평가사 甲은 乙주식회사가 소유한 △△동 1번지 소재 업무용빌딩과 △△동 1-1번지 나지상태의 토지에 대하여 재무보고목적의 감정평가를 진행하려 한다. 다음 물음에 답하시오. (30점) \par
1) 본건 감정평가의 기준가치는 무엇인지 그 개념에 관해 설명하고, 시장가치 기준원칙과의 관계에 관해 설명하시오. (10점) \par
2) 甲은 △△동 1번지 소재 업무용빌딩에 대하여 할인현금흐름분석법(discounted cash flow method)을 적용하려 한다. 이 때 적용할 할인율(discount rate)과 최종환원율(terminal capitalization rate)을 설명하고, 업무용 부동산시장의 경기변동과 관련하여 양자의 관계를 설명하시오. (15점) \par
3) △△동 1-1번지 토지에 대하여 공시지가기준법을 적용하여 시점수정, 지역요인 및 개별요인의 비교 과정을 거쳐 산정된 가액이 기준가치에 도달하지 못하였다고 가정할 경우 공시지가기준법에 따라 甲이 실무적으로 보정할 수 있는 방법에 관해 설명하시오. (5점) \par
\end{longquestion}
